% --- Archon physics formalization source begin ---
% archon:physics
% archon:covers PhyXMiniProblems/problem_phyx_mini_0065.lean
% archon:source-report reports/phyx_mini/problem_phyx_mini_0065.source.json

\chapter{Physics problem phyx\_mini\_0065}
\label{ch:PhyXMiniProblems_problem_phyx_mini_0065}

\paragraph{Problem source.}
\#\# Question

Find the focal length of the planoconvex polystyrene plastic lens in Figure

\#\# Answer choices

- A: A: 24cm

- B: B: 56cm

- C: C: 68cm

- D: D: 80cm

\#\# Figure caption (auxiliary; use the image as primary evidence)

The image depicts a diagram of a plano-convex lens. It includes:

- A plano-convex lens, shown with one flat side and one convex side, sitting on a horizontal line.
- A point marked on the line, at a distance of 40 cm from the lens, with an arrow indicating the measurement.
- Text below the lens that reads "Planoconvex lens."

\#\# Dataset metadata

Geometrical Optics; Physical Model Grounding Reasoning

\paragraph{Recorded answer/context.}
C: C: 68cm

\paragraph{Figure/image path.}
/root/proposal\_for\_physic/science-mango/phyx\_mini\_run/phyx\_data/test\_image/65.png

\paragraph{Formalization target.}
create a compiling Lean file with sorry bodies at `PhyXMiniProblems/problem\_phyx\_mini\_0065.lean`. The Lean declarations must preserve the physical quantities, dimensions or dimensional roles, figure labels, governing-law hypotheses, and final relation expressed by this problem.
Use Mathlib/Physlib names found through LeanExplore where available. If a physics API is missing, introduce faithful local abstractions rather than scalar placeholder aliases.

\begin{theorem}[Physics formalization target]
\label{thm:physics:phyx_mini_0065:target}
\lean{PhyXMiniProblems.ProblemPhyXMini0065.problem_phyx_mini_0065}
\uses{def:physics:phyx-mini-0065:phyxminiproblems-problemphyxmini0065-lengthincentimeters, def:physics:phyx-mini-0065:phyxminiproblems-problemphyxmini0065-planoconvexlenssetup, def:physics:phyx-mini-0065:phyxminiproblems-problemphyxmini0065-matchesproblemandfigure, def:physics:phyx-mini-0065:phyxminiproblems-problemphyxmini0065-usesstandardmaterialindices, def:physics:phyx-mini-0065:phyxminiproblems-problemphyxmini0065-hasphysicalopticalparameters, def:physics:phyx-mini-0065:phyxminiproblems-problemphyxmini0065-obeysparaxialplanoconvexlensmakerequation, def:physics:phyx-mini-0065:phyxminiproblems-problemphyxmini0065-matchesanswertonearestcentimeter}
The assigned autoformalize agent should translate this physics problem into Lean declarations in the covered file, with theorem and lemma proof bodies written as `by sorry`.
\end{theorem}
\begin{proof}
This is an autoformalization task, not a proof task. Produce statements that can later be proved by the physics prover without weakening the physical model.
\end{proof}

% --- Archon declaration topology begin ---
\section{Declaration topology}
\begin{definition}[Length Quantity]
  \label{def:physics:phyx-mini-0065:phyxminiproblems-problemphyxmini0065-lengthquantity}
  \lean{PhyXMiniProblems.ProblemPhyXMini0065.LengthQuantity}
  A signed physical length whose scalar readout changes coherently with units
\end{definition}
\begin{proof}
  This is the declared model definition of Length Quantity.
\end{proof}
\begin{definition}[centimeter Unit Choices]
  \label{def:physics:phyx-mini-0065:phyxminiproblems-problemphyxmini0065-centimeterunitchoices}
  \lean{PhyXMiniProblems.ProblemPhyXMini0065.centimeterUnitChoices}
  SI base units with centimeters selected as the length unit
\end{definition}
\begin{proof}
  This is the declared model definition of centimeter Unit Choices.
\end{proof}
\begin{definition}[length In Centimeters]
  \label{def:physics:phyx-mini-0065:phyxminiproblems-problemphyxmini0065-lengthincentimeters}
  \lean{PhyXMiniProblems.ProblemPhyXMini0065.lengthInCentimeters}
  \uses{def:physics:phyx-mini-0065:phyxminiproblems-problemphyxmini0065-lengthquantity, def:physics:phyx-mini-0065:phyxminiproblems-problemphyxmini0065-centimeterunitchoices}
  The scalar centimeter readout of a dimensionful length
\end{definition}
\begin{proof}
  This is the declared model definition of length In Centimeters.
\end{proof}
\begin{definition}[Optical Medium]
  \label{def:physics:phyx-mini-0065:phyxminiproblems-problemphyxmini0065-opticalmedium}
  \lean{PhyXMiniProblems.ProblemPhyXMini0065.OpticalMedium}
  The two homogeneous media in the optical model
\end{definition}
\begin{proof}
  This is the declared model definition of Optical Medium.
\end{proof}
\begin{definition}[Lens Face]
  \label{def:physics:phyx-mini-0065:phyxminiproblems-problemphyxmini0065-lensface}
  \lean{PhyXMiniProblems.ProblemPhyXMini0065.LensFace}
  The two physical faces of the lens
\end{definition}
\begin{proof}
  This is the declared model definition of Lens Face.
\end{proof}
\begin{definition}[Axis Side]
  \label{def:physics:phyx-mini-0065:phyxminiproblems-problemphyxmini0065-axisside}
  \lean{PhyXMiniProblems.ProblemPhyXMini0065.AxisSide}
  Left and right sides of the horizontal principal axis in the figure
\end{definition}
\begin{proof}
  This is the declared model definition of Axis Side.
\end{proof}
\begin{definition}[Radius Arrow Endpoint]
  \label{def:physics:phyx-mini-0065:phyxminiproblems-problemphyxmini0065-radiusarrowendpoint}
  \lean{PhyXMiniProblems.ProblemPhyXMini0065.RadiusArrowEndpoint}
  Geometrical roles of the two endpoints of the 40 cm arrow
\end{definition}
\begin{proof}
  This is the declared model definition of Radius Arrow Endpoint.
\end{proof}
\begin{definition}[Lens Surface Profile]
  \label{def:physics:phyx-mini-0065:phyxminiproblems-problemphyxmini0065-lenssurfaceprofile}
  \lean{PhyXMiniProblems.ProblemPhyXMini0065.LensSurfaceProfile}
  \uses{def:physics:phyx-mini-0065:phyxminiproblems-problemphyxmini0065-lengthquantity}
  A lens face is either planar or spherical with a physical radius magnitude
\end{definition}
\begin{proof}
  This is the declared model definition of Lens Surface Profile.
\end{proof}
\begin{definition}[Plano Convex Lens]
  \label{def:physics:phyx-mini-0065:phyxminiproblems-problemphyxmini0065-planoconvexlens}
  \lean{PhyXMiniProblems.ProblemPhyXMini0065.PlanoConvexLens}
  \uses{def:physics:phyx-mini-0065:phyxminiproblems-problemphyxmini0065-lengthquantity, def:physics:phyx-mini-0065:phyxminiproblems-problemphyxmini0065-opticalmedium, def:physics:phyx-mini-0065:phyxminiproblems-problemphyxmini0065-lensface, def:physics:phyx-mini-0065:phyxminiproblems-problemphyxmini0065-lenssurfaceprofile}
  The planoconvex lens and the physical quantity requested by the problem
\end{definition}
\begin{proof}
  This is the declared model definition of Plano Convex Lens.
\end{proof}
\begin{definition}[Radius Arrow Annotation]
  \label{def:physics:phyx-mini-0065:phyxminiproblems-problemphyxmini0065-radiusarrowannotation}
  \lean{PhyXMiniProblems.ProblemPhyXMini0065.RadiusArrowAnnotation}
  \uses{def:physics:phyx-mini-0065:phyxminiproblems-problemphyxmini0065-lengthquantity, def:physics:phyx-mini-0065:phyxminiproblems-problemphyxmini0065-radiusarrowendpoint}
  The dimensionful radius annotation printed in the primary figure
\end{definition}
\begin{proof}
  This is the declared model definition of Radius Arrow Annotation.
\end{proof}
\begin{definition}[Plano Convex Figure]
  \label{def:physics:phyx-mini-0065:phyxminiproblems-problemphyxmini0065-planoconvexfigure}
  \lean{PhyXMiniProblems.ProblemPhyXMini0065.PlanoConvexFigure}
  \uses{def:physics:phyx-mini-0065:phyxminiproblems-problemphyxmini0065-axisside, def:physics:phyx-mini-0065:phyxminiproblems-problemphyxmini0065-radiusarrowannotation}
  Qualitative geometry and quantitative annotation read from the figure
\end{definition}
\begin{proof}
  This is the declared model definition of Plano Convex Figure.
\end{proof}
\begin{definition}[Plano Convex Lens Setup]
  \label{def:physics:phyx-mini-0065:phyxminiproblems-problemphyxmini0065-planoconvexlenssetup}
  \lean{PhyXMiniProblems.ProblemPhyXMini0065.PlanoConvexLensSetup}
  \uses{def:physics:phyx-mini-0065:phyxminiproblems-problemphyxmini0065-planoconvexlens, def:physics:phyx-mini-0065:phyxminiproblems-problemphyxmini0065-planoconvexfigure}
  The lens together with its associated source figure
\end{definition}
\begin{proof}
  This is the declared model definition of Plano Convex Lens Setup.
\end{proof}
\begin{definition}[Matches Problem And Figure]
  \label{def:physics:phyx-mini-0065:phyxminiproblems-problemphyxmini0065-matchesproblemandfigure}
  \lean{PhyXMiniProblems.ProblemPhyXMini0065.MatchesProblemAndFigure}
  \uses{def:physics:phyx-mini-0065:phyxminiproblems-problemphyxmini0065-lengthincentimeters, def:physics:phyx-mini-0065:phyxminiproblems-problemphyxmini0065-planoconvexlenssetup}
  Problem and primary-figure readouts. The arrow endpoints record its geometric meaning, and its indicated length is identified with the lens's spherical radius. No focal-length value occurs in these data
\end{definition}
\begin{proof}
  This is the declared model definition of Matches Problem And Figure.
\end{proof}
\begin{definition}[Uses Standard Material Indices]
  \label{def:physics:phyx-mini-0065:phyxminiproblems-problemphyxmini0065-usesstandardmaterialindices}
  \lean{PhyXMiniProblems.ProblemPhyXMini0065.UsesStandardMaterialIndices}
  \uses{def:physics:phyx-mini-0065:phyxminiproblems-problemphyxmini0065-planoconvexlenssetup}
  The standard dimensionless refractive-index data used for the named media: ordinary air has index 1, and polystyrene plastic has index 1.59. These are material data, not a statement of the requested focal length
\end{definition}
\begin{proof}
  This is the declared model definition of Uses Standard Material Indices.
\end{proof}
\begin{definition}[Has Physical Optical Parameters]
  \label{def:physics:phyx-mini-0065:phyxminiproblems-problemphyxmini0065-hasphysicalopticalparameters}
  \lean{PhyXMiniProblems.ProblemPhyXMini0065.HasPhysicalOpticalParameters}
  \uses{def:physics:phyx-mini-0065:phyxminiproblems-problemphyxmini0065-lengthincentimeters, def:physics:phyx-mini-0065:phyxminiproblems-problemphyxmini0065-opticalmedium, def:physics:phyx-mini-0065:phyxminiproblems-problemphyxmini0065-planoconvexlenssetup}
  Positivity and optical-density conditions selecting the physical branch
\end{definition}
\begin{proof}
  This is the declared model definition of Has Physical Optical Parameters.
\end{proof}
\begin{definition}[Obeys Paraxial Plano Convex Lensmaker Equation]
  \label{def:physics:phyx-mini-0065:phyxminiproblems-problemphyxmini0065-obeysparaxialplanoconvexlensmakerequation}
  \lean{PhyXMiniProblems.ProblemPhyXMini0065.ObeysParaxialPlanoConvexLensmakerEquation}
  \uses{def:physics:phyx-mini-0065:phyxminiproblems-problemphyxmini0065-planoconvexlenssetup}
  The paraxial thin-lens lensmaker equation for one planar face and one convex spherical face, 1 / f = (n\_lens / n\_air - 1) / R. The equation is required in every unit choice, so it is dimensionally coherent. It is a governing physical law and contains no numerical value for the requested focal length
\end{definition}
\begin{proof}
  This is the declared model definition of Obeys Paraxial Plano Convex Lensmaker Equation.
\end{proof}
\begin{definition}[Answer Choice]
  \label{def:physics:phyx-mini-0065:phyxminiproblems-problemphyxmini0065-answerchoice}
  \lean{PhyXMiniProblems.ProblemPhyXMini0065.AnswerChoice}
  Labels of the four displayed multiple-choice answers
\end{definition}
\begin{proof}
  This is the declared model definition of Answer Choice.
\end{proof}
\begin{definition}[answer Focal Length In Centimeters]
  \label{def:physics:phyx-mini-0065:phyxminiproblems-problemphyxmini0065-answerfocallengthincentimeters}
  \lean{PhyXMiniProblems.ProblemPhyXMini0065.answerFocalLengthInCentimeters}
  \uses{def:physics:phyx-mini-0065:phyxminiproblems-problemphyxmini0065-answerchoice}
  The whole-centimeter focal-length readout printed beside each choice
\end{definition}
\begin{proof}
  This is the declared model definition of answer Focal Length In Centimeters.
\end{proof}
\begin{definition}[Matches Answer To Nearest Centimeter]
  \label{def:physics:phyx-mini-0065:phyxminiproblems-problemphyxmini0065-matchesanswertonearestcentimeter}
  \lean{PhyXMiniProblems.ProblemPhyXMini0065.MatchesAnswerToNearestCentimeter}
  \uses{def:physics:phyx-mini-0065:phyxminiproblems-problemphyxmini0065-answerchoice, def:physics:phyx-mini-0065:phyxminiproblems-problemphyxmini0065-answerfocallengthincentimeters}
  Agreement of an exact focal length with a displayed nearest-centimeter choice
\end{definition}
\begin{proof}
  This is the declared model definition of Matches Answer To Nearest Centimeter.
\end{proof}
\begin{lemma}[focal Length In Centimeters eq four thousand over fifty nine]
  \label{lem:physics:phyx-mini-0065:phyxminiproblems-problemphyxmini0065-focallengthincentimeters-eq-four-thousand-over-fifty-nine}
  \lean{PhyXMiniProblems.ProblemPhyXMini0065.focalLengthInCentimeters_eq_four_thousand_over_fifty_nine}
  \uses{def:physics:phyx-mini-0065:phyxminiproblems-problemphyxmini0065-lengthincentimeters, def:physics:phyx-mini-0065:phyxminiproblems-problemphyxmini0065-planoconvexlenssetup, def:physics:phyx-mini-0065:phyxminiproblems-problemphyxmini0065-matchesproblemandfigure, def:physics:phyx-mini-0065:phyxminiproblems-problemphyxmini0065-usesstandardmaterialindices, def:physics:phyx-mini-0065:phyxminiproblems-problemphyxmini0065-hasphysicalopticalparameters, def:physics:phyx-mini-0065:phyxminiproblems-problemphyxmini0065-obeysparaxialplanoconvexlensmakerequation}
  The radius, material indices, and lensmaker equation determine the exact thin-lens focal-length readout 4000/59 cm
\end{lemma}
\begin{proof}
  The conclusion follows from the governing relations and figure readouts named in the statement of focal Length In Centimeters eq four thousand over fifty nine.
\end{proof}
% --- Archon declaration topology end ---
% --- Archon physics formalization source end ---
